\documentclass[11pt]{article}
\usepackage[a4paper,margin=27mm]{geometry}
\usepackage{amsmath,amssymb,amsthm,hyperref}
\newtheorem{theorem}{Theorem}
\newtheorem{lemma}[theorem]{Lemma}
\newcommand{\PSH}{\operatorname{PSH}}
\newcommand{\Eone}{\mathcal E^1}
\newcommand{\OO}{\mathcal O}
\newcommand{\vol}{\operatorname{vol}}
\newcommand{\an}{\mathrm{an}}
\title{Section approximation of the minimal reference and\protect\\
general-output relative integrability in a big rational class}
\author{Unreviewed conditional research, continuation 196}
\date{6 October 2026}
\begin{document}
\maketitle
\begin{abstract}
We prove a direct section approximation of the normalized minimal
reference off one actual algebraic divisor. Together with algebraic
nullness and finite-atomic output moment estimates, this transfers
relative finite-energy test integrability to arbitrary weak limits of
model outputs. Every input is explicitly stated below. In particular
the relevant outer-energy construction and finite-atomic estimates are
unreviewed inputs, not established general equation theorems. We do not
prove weighted moving-product convergence, negligible-discrepancy
nullness, fixed-target strong integration, or the inverse equation.
\end{abstract}

\section{Actual domain and exact conditional inputs}
Let $X$ be smooth integral projective of dimension $n\geq1$ over an
algebraically closed trivially valued characteristic-zero field $k$.
Let $L$ be an actual big rational line bundle, $\theta=c_1(L)$,
$K=X^\an$, and $V=\vol(L)>0$. The carrier $K$ is compact Hausdorff;
supported semivaluations are retained. Write $v_0$ for the generic
trivial valuation and $R=P_\theta(0)$ for the normalized genuine
minimal reference. Fix the actual trivial reference metric convention
for section potentials: their local expressions are finite logarithmic
maxima and their changes of frame are those of the indicated bundle.
The same convention is used in every tensor product below.

The following statements are assumptions of this supporting paper.
They have been proposed in the complete partial foundations and the
finite-nef and finite-atomic research papers; they require independent
reconstruction before being used for an original-scope solution.
\begin{enumerate}
\item The genuine $\PSH(\theta)$ cone consists of the all-point
limits of one decreasing net of finite, possibly polar, FS potentials
on actual rational bundles $M_i$ whose numerical classes tend to
$\theta$, with a finite generic endpoint. It has genuine model
envelopes and decreasing closure. The reference $R$ is genuine,
$R(v_0)=0$, and every normalized genuine $u$ satisfies $u\leq R\leq0$.
Every normalized FS potential on the actual bundle $L$ lies below $R$.
The raw energy $E$ is constant-equivariant and decreasing-continuous;
$E(R)=0$. Write $B(u)=-E(u)$ for normalized $u\in\Eone(\theta)$.
\item Model-envelope outputs $\mu_D$ are target independent, positive
finite divisorial atomic measures of raw mass $V$. For every normalized
model envelope $u_D$ and normalized genuine $v\in\Eone(\theta)$ they
satisfy the actual relative moment bound
\begin{align}
 \int(R-v)\,d\mu_D&\leq F_n(B(v),B(u_D)),\label{atomic}\tag{A}\\
 F_n(a,b)&=C_n a^{\alpha_n}b^{1/2}
                  \max(a,b)^{1/2-\alpha_n},\qquad\alpha_n=2^{-n}.
 \nonumber
\end{align}
Here $F_n$ is extended continuously by zero when $a=0$ or $b=0$.
No probability normalization is made. This is the all-test conclusion
of the finite-atomic local squeeze, conditional on its own foundations.
\item For every normalized $u\in\Eone(\theta)$ there are decreasing
model-majorant envelopes $\widetilde u_i\downarrow u$ with
$E(\widetilde u_i)\downarrow E(u)$ and
$d_i=\widetilde u_i(v_0)\downarrow0$. Their normalized envelopes
$u_i=\widetilde u_i-d_i$ have $B(u_i)\to B(u)$, their outputs are
unchanged by this constant shift, and
$\mu_i=\mu_{D_i}\to\mu_u$ weakly. The normalized $u_i$ are not asserted
to decrease. The positive Radon measure $\mu_u$ has raw mass $V$.
\item For every actual nonzero algebraic section $s$ of a rational
bundle, $\mu_u(Z(s)^\an)=0$. This is the algebraic section-zero
nullness of the proposed outer-energy construction. It does not state
that nonalgebraic polar or negligible sets are null.
\end{enumerate}
These hypotheses do not include the conclusion about general outputs
proved below. Standard Kodaira decomposition and eventual generation
of an actual ample rational bundle are ordinary algebraic background.
The earlier finite-nef draft's displayed energy identity has a coefficient
orientation error: with factors $T_\xi^jT_0^{n-1-j}$ the weights are
$n-j$, not $j+1$. The newer finite-model reconstruction retains the
same dimension-only moment estimate with the corrected weights.
Input 2 assumes that estimate itself; the erroneous identity is not a premise.

\section{A full-section supremum off a Kodaira divisor}
Choose an actual ample rational bundle $A$ and rational $a>0$ so that
$L-aA=\OO(E_0)$ for an actual effective rational Cartier divisor $E_0$.
Write $W=q_{E_0}$ for its single-section potential. Thus $W(v_0)=0$,
$W\leq0$, and $W$ is finite exactly off $Z=E_0^\an$. It is an FS
potential on the actual bundle $L-aA$. The bundle identity, and not
merely an equality in numerical classes, is used here.

For every sufficiently divisible nonempty degree $m$, let
\[
 \phi_m=\frac1m\max_{s\in\mathcal B_m}\log|s|,
 \qquad \mathcal B_m\text{ a basis of }H^0(X,mL),
\]
in the stated trivial reference convention. Its value is independent
of the choice of basis: trivial valuations satisfy the ultrametric
inequality and all nonzero scalar coefficients have absolute value one.
These functions are continuous with values in $[-\infty,0]$ and
are genuine normalized FS potentials on $L$. The powers of sections
give $\phi_m\leq\phi_{m'}$ when $m\mid m'$. The degrees are directed
by divisibility, and cofinal degrees can include any finitely many
denominator or generation requirements.

\begin{lemma}\label{reference}
For every $x\in K\setminus Z$,
\[
                    R(x)=\sup_m\phi_m(x).             \tag{1}
\]
In particular $R$ is finite and continuous on the ordinary open set
$K\setminus Z$. No equality on $Z$ is asserted.
\end{lemma}
\begin{proof}
The inequality $\phi_m\leq R$ is part of the envelope assumption.
Take a defining FS net $h_i\downarrow R$ on actual bundles $M_i$.
Normalize each stage by $\widehat h_i=h_i-h_i(v_0)$. Its generic
value is zero, its finite maximum weights are nonpositive, and it is
nonpositive on the entire carrier. Although these normalized stages
need not decrease, $\widehat h_i(x)\to R(x)$ at every point, since
$h_i(v_0)\to0$.

Fix a rational $t\in(0,1)$ before choosing a late stage. The actual
rational bundle
\[
 B_{i,t}=L-(1-t)M_i-t(L-aA)
\]
is ample for all sufficiently late $i$: its numerical class tends to
$ta[A]$, and the ample cone is open. A sufficiently divisible power
of $B_{i,t}$ is globally generated. The full generated-system FS
potential on this bundle is zero everywhere. Indeed at the center of
every original semivaluation at least one generating section is a unit
in a local frame; its value is zero, and all other section values are
nonpositive. This includes semivaluations with support.

Clear the denominators of $t$, the bundles, and the FS stages and
tensor their actual sections. The finite FS expression
\[
                  f_{i,t}=(1-t)\widehat h_i+tW
\]
is consequently on the actual bundle $L$: the generated factor on
$B_{i,t}$ contributes zero. The maximum weights of $f_{i,t}$ are
nonpositive and its generic value is zero. Each tensor section belongs
to one actual $H^0(X,mL)$, so its weighted maximum is at most $\phi_m$.
Real FS weights cause no issue; only tensor exponents and actual bundle
denominators need to be integral after this clearing. Therefore
\[
             \sup_m\phi_m(x)\geq(1-t)\widehat h_i(x)+tW(x).
\]
For $x\notin Z$, $W(x)$ is finite. Also $R(x)$ is finite: tensor the
section of $E_0$ with one full generated power of $aA$ to obtain an
actual full-system lower bound $\phi_{m_0}(x)\geq W(x)>-\infty$.
Let $i$ tend along its defining tail, and then let rational $t\downarrow0$.
This proves (1) with the parameters in that order.

On $K\setminus Z$, all cofinal $\phi_m$ after $m_0$ are finite
continuous functions. Their supremum is lsc. The reference $R$ is usc
by genuineness; hence their common finite value is continuous there.
\end{proof}

The same proof does not give full-section minorants of every unbounded
$v$. Its envelope step uses only the zero obstacle for $R$. Nor does
it identify a section supremum with its usc regularization on $Z$.

\section{Lower semicontinuity and the general output bound}
\begin{lemma}\label{lsc}
For a normalized genuine $v\in\Eone(\theta)$ define the function
\[
 b_v(x)=
 \begin{cases}R(x)-v(x),&x\notin Z,\\0,&x\in Z.\end{cases}       \tag{2}
\]
Then $b_v:K\to[0,+\infty]$ is lsc. The zero on $Z$ is an auxiliary
extension, not a subtraction of infinite values.
\end{lemma}
\begin{proof}
The reference is finite continuous on $K\setminus Z$, and $v$ is usc,
so its difference there is extended lsc. It is nonnegative because
$v\leq R$. At a point of the closed set $Z$ the assigned value is
zero and all nearby values are nonnegative, which is precisely the
lower semicontinuity inequality at that point. Equivalently every
positive superlevel set of $b_v$ is open within the open complement
of $Z$, hence is open in $K$.
\end{proof}

For clarity, the extended Portmanteau inequality used next holds for
nets of positive Radon measures on compact Hausdorff spaces:
\[
         \int b\,d\nu\leq\liminf_i\int b\,d\nu_i
         \quad(b\geq0\text{ lsc},\ \nu_i\to\nu\text{ weakly}).
                                                               \tag{3}
\]
For bounded truncations, approximate the Radon integral from below by
continuous minorants. Finite maxima make these minorants directed;
weak convergence applies to each fixed continuous minorant. Then take
their supremum and let the truncation height increase. This proves
(3) without a countability, metrization, or sequential-compactness
assumption.

\begin{theorem}\label{general}
Under the four exact conditional inputs, every normalized genuine
$u,v\in\Eone(\theta)$ satisfies
\[
      \int(R-v)\,d\mu_u\leq F_n(B(v),B(u))<+\infty.       \tag{4}
\]
The relative difference is taken on the complement of the reference
pole set, which is $\mu_u$-null. In particular $\mu_u$ gives zero mass
to $\{v=-\infty\}\setminus Z$ for each such $v$.
\end{theorem}
\begin{proof}
The outputs $\mu_i$ in input 3 have finite divisorial support. No such
function-field valuation lies on $Z$: a nonzero equation of $E_0$
has a finite value at it. Thus
$\int b_v\,d\mu_i=\int(R-v)\,d\mu_i$. Apply (3), input 2, and the
continuity of $F_n$ to obtain
\[
 \int b_v\,d\mu_u
 \leq\liminf_i\int b_v\,d\mu_i
 \leq\lim_i F_n(B(v),B(u_i))=F_n(B(v),B(u)).
\]
Input 4 gives $\mu_u(Z)=0$. Lemma~\ref{reference} gives
$\{R=-\infty\}\subset Z$. Hence the left integral is exactly the
a.e. defined relative integral in (4), and no mass at a common pole
has been discarded without proof. On $\{v=-\infty\}\setminus Z$
the function $b_v$ is infinite, so its finite integral implies the
last assertion.
\end{proof}

\section{Exact scope and remaining work}
This is a conditional transfer of \emph{relative integrability}, with
an explicit uniform moment bound, from finite model outputs to the
proposed general outer-energy outputs. It retains every characteristic-zero
field, actual nonnef big rational bundle, supported semivaluation and
raw mass in the fixed domain. The proof does not replace that domain
by an abstract compact completion or a nef reference bundle.

It proves no absolute integrability of $R$: (4) controls only $R-v$.
It proves no nullness for all nonalgebraic pluripolar or negligible sets,
since no theorem representing all of them by finite-energy tests has
been supplied. It proves no weighted convergence for a changing
unbounded moving quotient. Its weak-measure step is a one-sided lsc
bound, not convergence of the integrals. In particular it cannot by
itself justify absolute $C(K)$ approximants with uniform $L^1$ error,
strong integration against each intrinsic target, variational attainment,
strict Borel comparison, or uniqueness of every equation solution.
Original P16.5 remains open; no accepted-paper count changes.

The section argument uses ordinary Kodaira decomposition, the open
ample cone and eventual global generation. The finite-atomic estimate
uses finite-nef Hodge and complete moving products through its separately
stated unreviewed foundations. The algebraic-nullness and weak-continuity
inputs use the separately reconstructed outer-energy derivative. None
of these is silently imported from a historical verdict or hash.
No active campaign output was consulted and no hot feedback was sent.

\begin{thebibliography}{9}
\small
\bibitem{BJ} S. Boucksom and M. Jonsson,
\emph{Global pluripotential theory over a trivially valued field},
\href{https://arxiv.org/pdf/1801.08229v4}{arXiv:1801.08229v4}.
The ample FS convention and nonnegligibility background do not themselves
give the nonample general-output theorem.
\bibitem{partial} Complete original-scope partial foundations,
model-envelope energy derivatives and algebraic section-zero nullness,
14-page research draft, 6 October 2026. These are inputs 1, 3 and 4,
not independently accepted premises in this paper.
\bibitem{atomic} Finite-nef moment and mixed-moving estimates;
finite-divisorial-output local squeezing and all finite-energy test
extension, two 4-page conditional research drafts, 6 October 2026.
Their precise all-test conclusion is input 2.
\end{thebibliography}
\end{document}
